\documentclass[12pt]{article}
\usepackage[utf8]{inputenc}
\usepackage[T1]{fontenc}
\usepackage[spanish]{babel}
\usepackage{amsmath, amssymb, amsfonts}
\usepackage{fancyhdr}
\usepackage{xcolor}
\usepackage[margin=3cm]{geometry}
\usepackage{graphicx}
\usepackage{float}
\usepackage{multicol}
\usepackage{enumitem}
\usepackage{tcolorbox}
\usepackage{listings}

% Configuración de listings para Python
\definecolor{codegreen}{rgb}{0,0.6,0}
\definecolor{codegray}{rgb}{0.5,0.5,0.5}
\definecolor{codepurple}{rgb}{0.58,0,0.82}
\definecolor{backcolour}{rgb}{0.95,0.95,0.92}

\lstdefinestyle{mystyle}{
    backgroundcolor=\color{backcolour},   
    commentstyle=\color{codegreen},
    keywordstyle=\color{magenta},
    numberstyle=\tiny\color{codegray},
    stringstyle=\color{codepurple},
    basicstyle=\ttfamily\small,
    breakatwhitespace=false,         
    breaklines=true,                 
    captionpos=b,                    
    keepspaces=true,                 
    numbers=left,                    
    numbersep=5pt,                  
    showspaces=false,                
    showstringspaces=false,
    showtabs=false,                  
    tabsize=2,
    language=Python
}

\lstset{style=mystyle}

% Encabezado
\pagestyle{fancy}
\fancyhf{}
\fancyhead[L]{Escuela de Ingeniería \\ Universidad de O'Higgins}
\fancyhead[R]{Programación I \\ 2026-01}
\fancyfoot[C]{\thepage}

% Título comprimido
\makeatletter
\renewcommand{\maketitle}{%
  \begin{center}
    {\LARGE \@title\par}
    \vskip 0.3em
    {\large \@author\par}
    \vskip 0.3em
  \end{center}
  \vskip -0.2em  % Espacio después del título
}
\makeatother

\title{Guía de Ayudantía: Funciones en Python}
\author{Felipe Llach R.}

\begin{document}
\maketitle

\section{Introducción}
Las funciones son bloques de código reutilizables que realizan una tarea específica. Nos permiten organizar el código, evitar la repetición y facilitar la resolución de problemas complejos dividiéndolos en partes más pequeñas.

\section{Definición de una Función}
Para definir una función en Python se utiliza la palabra reservada \texttt{def}, seguida del nombre de la función y paréntesis.

\begin{tcolorbox}[title=Estructura básica]
\begin{lstlisting}
def nombre_de_la_funcion(parametros):
    # Cuerpo de la funcion
    instrucciones
    return resultado # Opcional
\end{lstlisting}
\end{tcolorbox}

\section{Parámetros y Argumentos}
Los \textbf{parámetros} son las variables que se listan entre los paréntesis en la definición de la función. Los \textbf{argumentos} son los valores que se envían a la función cuando esta es llamada.

\begin{lstlisting}
def saludar(nombre):
    print(f"Hola, {nombre}!")

saludar("Diego") # "Diego" es el argumento
\end{lstlisting}

\section{Retorno de Valores}
Una función puede devolver un valor al programa principal utilizando la sentencia \texttt{return}. Una vez que se ejecuta un \texttt{return}, la función termina inmediatamente.

\begin{lstlisting}
def sumar(a, b):
    return a + b

resultado = sumar(5, 3)
print(resultado) # Imprime 8
\end{lstlisting}

\section{Ejercicios Propuestos}

\begin{enumerate}
    \item \textbf{Calculadora de Área:} Escriba una función llamada \texttt{area\_triangulo(base, altura)} que reciba la base y la altura de un triángulo y retorne su área.
    
    \item \textbf{Verificador de Paridad:} Escriba una función \texttt{es\_par(numero)} que reciba un número entero y retorne \texttt{True} si el número es par y \texttt{False} si es impar.
    
    \item \textbf{Conversor de Temperatura:} Cree una función que convierta grados Celsius a Fahrenheit. La fórmula es: $F = C \times 1.8 + 32$.
    
    \item \textbf{Máximo de tres:} Escriba una función \texttt{max\_de\_tres(a, b, c) } que reciba tres números y retorne el mayor de ellos sin usar la función \texttt{max() } de Python.
\end{enumerate}

\end{document}
